\section{Aula 9}

Nessa aula continuamos o estudo de percolação em $\Z^2$ com o objetivo de chegar
no Teorema de Kesten.
\begin{theorem}
	$p_T(\Z^2) = p_H(\Z^2) = 1/2$ e $\Theta(1/2) = 0$.
\end{theorem}

Vimos o método de Russo-Seymour-Welsh e lemas de cruzamento de retângulos
e como eles implicam comportamento sub-crítico e super-crítico. Como não
serei nem um pouco capaz de reproduzir todos os desenhos feitos em
aula (foi só isso basicamente),
eu vou por aqui lemas equivalentes que estão no Bela.

Precisamos de uma observação que já foi vista antes, é uma
consequência fácil da dualidade
de $\Z^2$. Novamente, dado um retângulo $R \subset \Z^2$ com lados
$k$ por $l$, definimos $H(R)$
o evento que existe um caminho aberto cruzando $R$ horizontalmente e
$V(R)$ o mesmo verticalmente.
Notando que cruzamos horizontalmente $R$ se e somente se não
cruzamos verticalmente no dual (e vice versa), temos a seguinte observação.

\begin{observation}
	\begin{enumerate}
		\item Se $R$ e $R'$ são retângulos $k \times l-1$ e $k-1 \times l$
			respectivamente, então
			\[
				\Pr_p(H(R)) + \Pr_{1-p}(V(R')) = 1.
			\]
		\item Se $R$ é um retângulo $n+1$ por $n$, então
			\[\Pr_{1/2}(V(R)) = 1/2.\]

		\item se $Q$ é um quadrado $n$ por $n$, então
			\[\Pr_{1/2}(V(Q)) = \Pr_{1/2}(H(Q)) \geq 1/2\]
	\end{enumerate}
\end{observation}

Para facilitar nossa vida, em vez de fazer a prova do Augusto do lema
de Russo-Seymour-Welsh
que envolve separação de casos nas formas que estamos cruzando um
retângulo, vamos fazer a prova
do Bela.

\begin{lemma}
	Seja $R = [m]\times[2n]$, $m \geq n$, um retângulo. Seja $X(R)$ o
	evento que existem caminhos
	abertos $P_1$ e $P_2$ tais que $P_1$ cruza verticalmente o quadrado
	$Q = [n]\times[n]$
	e $P_2$ encosta em $P_1$ em $Q$ e toca o lado direito de $R$. Então
	\[
		\Pr_p(X(R)) \geq \Pr_p(H(R))\Pr_p(V(Q))/2.
	\]
\end{lemma}

\begin{remark}
	Ver Figura 5 do cap. 2 do Bela. O Augusto fez esse lema durante a sua prova.
\end{remark}

\begin{proof}
	A prova disso é bem bonitinha. Suponha que $V(Q)$ acontece, então
	existe um caminho $P_1$ mais
	a esquerda que separa $Q$ num pedaço esquerdo e num pedaço direito.
	Crucialmente, fazendo
	uma exploração boba, podemos revelar $P_1$ independentemente das
	arestas do lado direito
	de $Q$. Podemos refletir $P_1$ para o quadrado $[n]\times [n+1,2n]$
	em um outro caminho $P_1'$.
	Agora suponha que $H(R)$ acontece, como não revelamos nada a direita
	de $P_0$, o caminho
	$P_2$ que cruza $R$ horizontalmente tem a mesma chance de bater em
	$P_1$ quanto $P_1'$. Usando
	Harris, segue que $\Pr(X(R)) \geq \Pr(V(Q))\Pr(H(R))/2$.
\end{proof}

Esse lema permite, usando a observação, calcular bons lowerbounds na
probabilidade de cruzar
retângulos. Por exemplo, definindo $h_{1/2}(m,n) = \Pr(H([m] \times
[n]))$, temos

\begin{corollary}
	Para todo $n \geq 1$, $h(3n,2n) \geq 2^{-7}$.
\end{corollary}
\begin{proof}
	Seja $Q_1 = [2n]\times[2n]$ e $Q_2 = [n + 1, 3n] \times [2n]$,
	quadrados de lado $2n$
	inseridos no retângulo $R = [3n]\times[2n]$. Seja $S = [n+1,2n]
	\times [n]$ a metade de
	baixo da interseção de $Q_1$ e $Q_2$. Segue que
	\[
		H(R) \subset X(Q_1) \cap X(Q_2) \cap H(S).
	\]
	Portanto, por Harris e pela observação,
	\[
		\Pr(H(R)) \geq 2^{-7}.
	\]
\end{proof}

Isso implica o que vimos no final da aula. Se $p \leq 1/2$, então, no
dual, retângulos arbitrariamente
grandes são  cruzados no lado grande com probabilidade positiva
constante. Em particular,
considerando uma bola enorme $B_N$ ao redor do $0$, podemos contornar
ela com cercas (feitas desses
retângulos) disjuntas que estão fechadas com probabilidade positiva.
Segue imediatamente que
eventualmente alguma dessas cercas está fechada e temos decaimento
polinomial. Fazendo
as continhas e os desenhos (Figura 8 Cap 2 Bela), e sendo $N = C^k$
para algum $k$, vemos
\[
	\Pr(0 \leftrightarrow \partial B_{N}) \leq (1-\eps)^k = n^{-c}.
\]
O que prova $\Theta(1/2) = 0$. Na próxima aula veremos que para $p >
1/2$, conseguimos
fazer componentes arbitrariamente grandes também, o que finaliza o
teorema de Kesten.
